\documentclass[12pt, addpoints]{exam}
\usepackage[utf8]{inputenc}
\usepackage[T1]{fontenc}
\usepackage[brazilian]{babel}
\usepackage{amsmath, amssymb}
\usepackage{graphicx}
\usepackage{tikz}
\usepackage{pgfplots}
\pgfplotsset{compat=1.18}
\usepackage{booktabs}
\usepackage{xcolor}
\usepackage{mdframed}
\usepackage[top=1.5cm, bottom=1.5cm, left=2cm, right=2cm, headheight=14pt, headsep=8pt, footskip=10pt, includehead, includefoot]{geometry}

\marksnotpoints
\pointsinrightmargin
\marginpointname{ pts}

% Configuracao de Cabecalho e Rodape
\pagestyle{headandfoot}
\runningheader{Álgebra Linear}{}{Página \thepage\ de \numpages}
\footer{}{}{}



\begin{document}

% -----------------------------------------------------------
% Cabecalho
% -----------------------------------------------------------
\begin{center}
    \textsc{\textbf{Universidade Federal do Pará} \\
    CAMTUC -- FEE: Álgebra Linear} -- Prof.\ Dr.\ Raphael Teixeira \\
    \vspace{0.1cm}
    Avaliação 3: Espaços com Produto Interno \hspace{0.6cm}-\hspace{0.6cm} \rule{2.5cm}{0.1pt}/2026.
\end{center}
\vspace{0.3cm}
\noindent
\textbf{Aluno(a):} \rule{8cm}{0.1pt} \quad \textbf{Matrícula:} \rule{3.5cm}{0.1pt}
\vspace{0.0cm}
% -----------------------------------------------------------
\begin{questions}

% ============================================================
%  QUESTAO 1 -- Comprimento, Vetor Unitario e Distancia (5.1)
% ============================================================
\question
Considere os vetores $u=(1,\,2,\,-2,\,4)$ e $v=(3,\,-1,\,0,\,2)$, em $\mathbb{R}^4$.

\textbf{(a, 1.0 pt)} Calcule o comprimento (norma) $\|u\|$.

\textbf{(b, 1.0 pt)} Determine o vetor unitário na direção de $u$.

\textbf{(c, 1.5 pt)} Calcule a distância $d(u,v)$ entre os vetores $u$ e $v$.

% ============================================================
%  QUESTAO 2 -- Produto Escalar e Projecao Ortogonal (5.1)
% ============================================================
\question
Considere os vetores $u=(2,\,-1,\,3)$ e $v=(1,\,2,\,-2)$, em $\mathbb{R}^3$.

\textbf{(a, 1.0 pt)} Calcule o produto escalar $u\cdot v$.

\textbf{(b, 1.0 pt)} Calcule $v\cdot v$.

\textbf{(c, 1.5 pt)} Determine a projeção ortogonal de $u$ em $v$, isto é, $\text{proj}_v\,u = \dfrac{u\cdot v}{v\cdot v}\,v$.

% ============================================================
%  QUESTAO 3 -- Processo de Gram-Schmidt: Ortogonalizacao e Ortonormalizacao (5.3)
% ============================================================
\question
Considere a base $B=\{v_1,\,v_2,\,v_3\}$ de $\mathbb{R}^3$, dada por
\[
  v_1=(1,\,1,\,1), \qquad v_2=(0,\,1,\,1), \qquad v_3=(0,\,0,\,1).
\]

\textbf{(a, 1.5 pt)} Aplique o processo de ortogonalização de Gram-Schmidt para obter uma base ortogonal $B'=\{w_1,\,w_2,\,w_3\}$ de $\mathbb{R}^3$.

\textbf{(b, 1.5 pt)} A partir da base ortogonal $B'$ obtida no item anterior, determine a base ortonormal $B''=\{u_1,\,u_2,\,u_3\}$, normalizando cada vetor $w_i$.

\end{questions}
\end{document}
